\chapterimage{images/18-bekstvo-kroz-svetove.jpg}

\chapter{Bekstvo kroz svetove}

\editorialnote{Putovanje kao kičma radnje; zadržati samo svetove i epizode koje menjaju odnos likova ili razumevanje većeg sukoba.}

Konzulov brod napustio je Hiperion pre nego što je Paks uspeo da zatvori čitav sistem.

Za Rola je to bio prvi odlazak sa sveta na kojem je rođen.

Za Aneu povratak u Galaksiju koju više nije poznavala.

Kada je ušla u Sfingu, Hegemonija je tek bila mrtva.

Kada je izašla, prošla su gotovo dva i po veka.

\bigskip

Martin Silenus im je rekao da krenu prema starim svetovima Mreže.

Nije im dao preciznu kartu.

Samo ideju.

Reka Tetis.

\bigskip

Rol je mislio da je starac pred smrt konačno izgubio razum.

Tetis je nekada bila vodeni put koji je prolazio kroz desetine svetova.

Reka bi na jednoj planeti ušla u dalekobacački portal, a nekoliko metara dalje nastavila pod drugim suncem.

Jahte su mogle danima ploviti kroz različite planetarne sisteme a da putnici gotovo i ne primete prelaze.

Padom Mreže svi ti portali su umrli.

Jedna velika međuzvezdana reka raspala se na desetine običnih reka.

\bigskip

Bar je tako čovečanstvo verovalo.

\bigskip

Prvo su krenuli prema Parvatiju.

Bio je to zabačen svet koji nikada nije predstavljao važan deo stare Mreže.

Putovanje je trajalo nekoliko dana njihovog vremena, dok su na Hiperionu prošli meseci.

Rol je prvi put osetio pravo značenje vremenskog duga.

Nije više mogao jednostavno da se vrati kući.

Čak i ako jednog dana to učini, svet koji će pronaći neće biti onaj koji je napustio.

\bigskip

Anea je takve stvari prihvatala lakše.

Već je preskočila skoro dva i po veka.

Nekoliko meseci nije joj mnogo značilo.

\bigskip

Tokom putovanja Rol je pokušavao da izvuče iz nje plan.

``Kuda zapravo idemo?''

``Prema starim svetovima.''

``To nije odgovor.''

``Jeste. Samo nije naročito dobar.''

``Tražimo Proterane?''

``Možda.''

``Onog arhitektu iz tvojih snova?''

``Da.''

``Na kom svetu?''

``Ne znam.''

Rol bi obično tada odustao.

Anea bi se vratila knjizi.

\bigskip

Nedelja provedena u brodu promenila je njihov odnos više nego dramatično bekstvo sa Hiperiona.

Jeli su zajedno.

Igrali šah.

Go.

Poker.

Anea je čitala gotovo sve što bi pronašla.

Rol je otkrio da dvanaestogodišnja devojčica ume pažljivije da sluša nego većina odraslih koje je poznavao.

Ona je zauzvrat izvukla iz njega priče o pastirskom životu, vojsci, baržama, kockarnicama i Izi.

\bigskip

A. Betik je uglavnom posmatrao.

Android je bio mnogo stariji od oboje.

Bio je biofakturisan pre skoro sedam vekova.

Služio je kraljevima.

Došao na Hiperion pre Tužnog Kralja Bilija.

Poznavao svet koji je za Rola postojao samo u pričama.

A sada je ponovo putovao njime.

\bigskip

Na Parvatiju ih je čekao Paks.

Ne slučajno.

Otac Kapetan de Soja predvideo je njihov pravac i stigao mnogo ranije pomoću kurirskog broda čija posada je mogla da podnese ubrzanja koja bi živog čoveka pretvorila u kašu.

Rešenje Paksa bilo je jednostavno.

Posada bi umrla tokom putovanja.

Kruciforme bi je po dolasku ponovo sastavile.

\bigskip

Rol prvi put nije razumeo koliko velika prednost to daje njihovim progoniteljima.

Anea jeste.

Zato nije nameravala da ostane na Parvatiju.

Njihov pravi cilj bio je Renesansa Vektor.

Stari svet Mreže.

Svet koji je nekada imao dalekobacače.

\bigskip

Paks ih je gotovo presreo.

Konzulov brod nije bio ratna letelica.

Njegova zaštitna polja nisu mogla da zaustave ozbiljno vojno oružje.

Ako ih flota opkoli, neće moći da se bore.

Anea je zato predložila plan koji je čak i Rolu delovao suludo.

Iskoristiće sposobnost broda da menja oplatu.

Ako vojnici Paksa uđu unutra, brod može otvoriti čitave palube prema vakuumu.

\bigskip

``A mi?'' pitao je Rol.

``Bićemo zaštićeni poljima.''

``Nadaš se.''

``Da.''

Rol je pogledao A. Betika.

``Ona ima dvanaest godina.''

``Svestan sam toga, ser.''

\bigskip

Plan je uspeo dovoljno dobro da pobegnu.

Ali potera nije prestala.

De Soja je sada znao kojim brodom putuju.

I znao je da idu prema starim portalima Mreže.

\bigskip

Na Renesansi Vektor dogodilo se nešto što nije trebalo da bude moguće.

Stari dalekobacač se uključio.

\bigskip

Posle više od dva veka mrtvi luk iz vremena Hegemonije ponovo je zasijao.

Nije bilo TehnoCentra koji bi ga aktivirao.

Nije bilo Mreže.

Nije bilo razloga da radi.

\bigskip

Anea se nije iznenadila koliko Rol.

``Proći će.''

``Kako znaš?''

``Sanjala sam.''

``Naravno.''

\bigskip

Paks ih je sustigao pre nego što su uspeli mirno da priđu portalu.

Brod je pogođen.

Pao je u reku.

Rol je mislio da su gotovi.

Ali brod se nije zaustavio.

Zaštitna polja su ih učvrstila dok je letelica udarala o dno, orala kroz blato i gurala se prema starom luku.

Paks je otvorio vatru.

Voda se pretvarala u paru.

\bigskip

A onda je portal proradio.

Pod vodom.

\bigskip

Konzulov brod prošao je kroz njega nekoliko metara ispod površine reke.

U jednom trenutku bili su na Renesansi Vektor.

U sledećem pod drugim nebom.

\bigskip

Paks nije mogao za njima.

Kada je Rol pitao Aneu zašto, ona je odgovorila kao da govori o živom biću.

``Portal nas je pustio.''

``Portal je mašina.''

``Ne sasvim.''

``A zašto neće pustiti njih?''

``Zato što nisu sa mnom.''

\bigskip

Rol nije znao šta da uradi sa tim odgovorom.

Ali činjenica je ostala.

Dalekobacači koje je čovečanstvo smatralo mrtvim počeli su da se otvaraju samo pred Aneom.

\bigskip

To je promenilo čitavo putovanje.

Više nisu morali da prelaze svetlosne godine Hokingovim pogonom.

Mogli su da prate ostatke Tetisa.

Od jednog mrtvog luka do drugog.

Od sveta do sveta.

\bigskip

Nisu znali ko otvara put.

Anea je tvrdila da to nije TehnoCentar.

``Srž nema kontrolu nad ovim portalima?'' pitao je Rol.

``Ne.''

``Onda ko ima?''

``Ne znam.''

To je bio jedan od retkih trenutaka kada joj je odmah poverovao.

Izgledala je jednako zbunjeno kao on.

\bigskip

Prvi od svetova na novoj Tetis bio je Mare Infinitus.

Ogroman okean pod neobičnim nebom.

Portal ih je izbacio u more.

Konzulov brod bio je oštećen.

Umesto udobnog putovanja kroz svemir, njih troje morali su da nastave vodom.

Napravili su splav.

Poneli ono što su mogli.

Hoking-prostirka služila im je kao izviđačko sredstvo i transport.

\bigskip

Za Rola je to u jednom smislu bio povratak starom životu.

Ponovo voda.

Struja.

Vetar.

Oprema.

Noćna straža.

Samo što više nije vodio baržu po Hiperionu.

Plovio je mrtvom međuzvezdanom rekom sa detetom koje Paks želi više od svega.

\bigskip

Na Mare Infinitusu prvi put su naišli na ljude.

I na Paks.

Svet je bio pod njegovom kontrolom.

Stanovništvo je živelo od okeana, dok su vojne jedinice nadgledale luke i ogromne ribarske platforme.

Postojao je i otpor.

Ljudi koji nisu prihvatili novi poredak živeli su izvan zvaničnog sistema.

\bigskip

Rol se sukobio sa vojnicima.

Nije želeo da ubija više nego što mora.

Čak i kada ga je jedan od njihovih ljudi pokušao uhvatiti, Rol se vratio da spase ranjenog neprijatelja od utapanja.

Taj postupak skoro ga je koštao života.

\bigskip

Bio je teško ranjen.

Anea i A. Betik izvukli su ga.

Nastavili su niz Tetis dok je Rol ležao sa medkompom priključenim na telo.

Rane su bile ozbiljne.

Još opasnija bila je infekcija koju je pokupio u okeanu.

\bigskip

Anea ga je negovala.

Menjala zavoje.

Dopunjavala infuziju.

Dala mu lekove kada više nije mogao da podnese bol.

Rol je pokušavao da se ponaša kao da sve drži pod kontrolom.

Nije.

\bigskip

``Nemoj da umreš'', rekla mu je.

``To mi je i plan.''

``Nije naročito dobar.''

``Učim od tebe.''

\bigskip

Portal ih je sledeći put odveo tamo gde niko nije očekivao.

Na Hebron.

\bigskip

Hebron nikada nije bio prava stanica reke Tetis.

U vreme Hegemonije dozvoljavao je samo jedan dalekobacač, u Novom Jerusalimu.

A ipak ih je portal doveo upravo tamo.

Još jedan dokaz da njihovo putovanje nije pratilo staru mrežu po pravilima koja je nekada imala.

\bigskip

Hebron ih je dočekao tišinom.

\bigskip

Novi Jerusalim bio je prazan.

Ne razoren.

Ne spaljen.

Prazan.

\bigskip

U gradu je nekada živelo više od milion ljudi.

Sada na ulicama nije bilo nikoga.

Nije bilo vojske.

Nije bilo Proteranih.

Nije bilo gomila mrtvih.

Ljudi su jednostavno nestali.

\bigskip

Za Rola je to bilo strašnije od ruševina.

Ruševina govori šta se dogodilo.

Prazan grad postavlja pitanje.

\bigskip

Nisu imali vremena da ga rešavaju.

Rol je umirao.

U Sinajskom medicinskom centru pronašli su funkcionalne automatske sisteme.

Autohirurg ga je preuzeo.

Tokom nekoliko dana izvršio je više operacija.

Mikroorganizam sa Mare Infinitusa već se širio kroz njegovo telo.

Da su stigli samo dan kasnije, verovatno ga ne bi bilo moguće spasiti.

\bigskip

Rol je proveo skoro dve nedelje na Hebronu oporavljajući se.

Anea i A. Betik istraživali su grad.

Nisu pronašli razumno objašnjenje za nestanak stanovništva.

\bigskip

Samo još jedno pitanje.

\bigskip

Za Rola je ipak nešto postalo jasnije.

Anea ga više nije trebala samo kao telohranitelja.

Dok je on bio bespomoćan, ona je čuvala njega.

Martinova naredba da je štiti polako je gubila prvobitno značenje.

Više nisu bili čuvar i teret.

Postajali su grupa.

\bigskip

Tetis ih je zatim odvela u gotovo suprotan svet.

Sol Drakoni Septem.

\bigskip

Led.

Ogromna gravitacija.

Smrznuta atmosfera.

Gradovi zakopani ispod slojeva leda.

\bigskip

Reka je tekla kroz ledene tunele.

U jednom trenutku put im je bio potpuno zatvoren.

Rol je mislio da je njihova sreća konačno nestala.

Anea je drhtala čak i pod termalnim ćebetom.

A. Betik se pod visokom gravitacijom jedva kretao.

\bigskip

Ipak su pronašli tragove života.

Čitčatuke.

Ljude prilagođene tom smrznutom svetu.

I starog sveštenika po imenu Otac Glaukus.

\bigskip

Glaukus nije pripadao Crkvi kakvu je Rol poznavao.

Bio je prognan.

Pripadao je jednom od tejarovskih redova koje je novi Paks smatrao opasnim.

Decenijama je živeo među domorocima.

Nije pokušavao da ih pretvori u podanike Crkve.

Više ga je zanimalo da ih razume.

\bigskip

Prvi put tokom putovanja Rol je upoznao sveštenika koji nije pokušavao nikoga da kontroliše.

\bigskip

Glaukus je sa Aneom dugo razgovarao o svesti.

O evoluciji.

O ideji da čovečanstvo nije završena forma.

Da se svesna bića možda kreću prema nečemu većem.

\bigskip

Rol je bio nepoverljiv kada je razgovor stigao do TehnoCentra.

Za njega su AI bile neprijatelj iz istorije.

Glaukus nije tako lako prihvatao podelu.

Ako je nešto samosvesno, govorio je, onda pitanje njegove moralnosti ne može biti rešeno samo time od čega je napravljeno.

Mašina može da izabere.

Čovek takođe.

I jedno i drugo može izabrati pogrešno.

\bigskip

Aneu je taj razgovor zanimao mnogo više nego Rola.

``Empatija'', rekla je.

Glaukus je klimnuo.

Upravo je to smatrao ključem.

Ne inteligenciju samu po sebi.

Sposobnost da drugo biće ne posmatraš samo kao sredstvo.

\bigskip

Rol se setio svega što je čuo o Ultimativnoj Inteligenciji TehnoCentra.

Čist razum.

Predviđanje.

Kontrola.

Sada je prvi put počeo da razume zašto se stalno vraća ista reč.

Empatija.

\bigskip

Kada su napustili Sol Drakoni Septem, Rol je mislio da je Glaukus ostao bezbedan među svojim ljudima.

Anea je kasnije znala da nije.

\bigskip

Na sledećem delu puta razbolela se.

Imala je groznicu i trenutke u kojima je govorila kao da je istovremeno na dva mesta.

Jednom je zgrabila Rola za ruku.

``Voliš li me?''

Rol je ostao zatečen.

``Naravno da mi je stalo do tebe.''

Anea ga je nekoliko trenutaka gledala kao nepoznatog čoveka.

Onda joj se izraz promenio.

``Izvini.''

``Za šta?''

``Pogrešila sam vreme.''

\bigskip

Rol nije razumeo.

Anea je zatvorila oči.

``Zaboravila sam ko smo sada.''

\bigskip

Nije želela dalje da objašnjava.

Ali rekla im je nešto drugo.

Otac Glaukus je mrtav.

I Čitčatuke koji su bili sa njim.

\bigskip

``Šrajk?'' pitao je Rol.

``Ne.''

Anea je bila sigurna.

``Nešto drugo.''

\bigskip

Tada im je prvi put jasno rekla ono čega se najviše plašila.

Paks nije pravi neprijatelj.

Paks ih progoni zato što ga neko usmerava.

\bigskip

TehnoCentar nije nestao.

\bigskip

``Srž je živa'', rekla je.

Rol je odmahnuo glavom.

``Mreža je uništena.''

``Mreža jeste.''

``Bez dalekobacača nije imala gde da živi.''

``Našla je način.''

``Gde?''

Anea ga je pogledala.

``U Paksu.''

\bigskip

Rol je želeo da joj kaže da to nema smisla.

Paks je zabranio AI.

Crkva je vekovima učila da je TehnoCentar izdao čovečanstvo.

\bigskip

Upravo zato je odgovor bio toliko strašan.

Ako je Anea u pravu, onda je institucija koja tvrdi da je spasla ljude od Srži postala njen novi zaklon.

\bigskip

Kako i zašto, još nije mogla potpuno da objasni.

Ali bila je sigurna u jedno.

Ono što ih prati nije samo vojska Crkve.

Nešto iz TehnoCentra želi da je zaustavi.

\bigskip

Put ih je dalje odveo na Kom-Rijad.

Pustinjski svet.

Crveni kamen.

Kanali za navodnjavanje.

Sela.

Gradovi.

I ponovo gotovo potpuna tišina.

\bigskip

Prizor je podsećao na Hebron.

Ostavljene kuće.

Napuštena polja.

Mesta koja bi trebalo da budu puna ljudi.

\bigskip

Aneu je tamo pogodio još jedan talas bola.

Ovoga puta Rol je odmah znao da nije u pitanju obična bolest.

Devojčica je videla nešto.

Ili se nečega setila.

Ne iz prošlosti.

Iz budućnosti.

\bigskip

``Ne razumem kako možeš da se sećaš nečega što se još nije dogodilo.''

``Ni ja.''

``Ali sigurna si?''

``Da.''

\bigskip

Anea je sada znala da je Glaukus mrtav zato što će tu činjenicu jednog dana saznati.

Za nju su se povremeno delovi budućeg života prelivali unazad.

Nije mogla da ih kontroliše.

Nije mogla pouzdano da razlikuje mogućnost od događaja koji će se stvarno zbiti.

\bigskip

To ju je plašilo.

Rol je prvi put video koliko je njeno navodno znanje zapravo teret.

\bigskip

Šrajk se pojavio na Kom-Rijadu.

Nije napao.

Samo je stajao kraj njihovog puta.

\bigskip

Rol je podigao oružje.

Anea mu je spustila ruku.

``Neće pomoći.''

``Je li na našoj strani?''

``Ne.''

``Na strani Paksa?''

``Ne znam.''

``TehnoCentra?''

Anea je odmahnula glavom.

``Ne znam ko ga je ovog puta poslao.''

\bigskip

To ``ovog puta'' Rolu se nije dopalo.

``Hoćeš da kažeš da može imati različite gospodare?''

``Možda.''

``Ili?''

``Možda više nema gospodara.''

\bigskip

Šrajk ih je posmatrao.

Zatim ih je pustio da prođu.

\bigskip

Rol je tek tada u potpunosti prihvatio da priče iz Speva više nisu istorija.

Isti monstrum koji je stajao pred Solom Vejntraubom.

Koji se borio sa Kasadom.

Koji je odneo Martina na Drvo bola.

Sada je bio deo Aneinog puta.

\bigskip

Te noći Anea je konačno rekla ono što je mesecima izbegavala.

``Ne juri nas Šrajk.''

``Dobro.''

``Ne juri nas ni Paks.''

Rol se nasmejao.

``Imamo dovoljno dokaza za suprotno.''

``Paks nas prati. Ali ne zato što sam njemu neprijatelj.''

``Nego?''

``Zato što Srž kaže da sam pretnja.''

\bigskip

Objasnila mu je da se TehnoCentar posle Pada povukao, ali nije nestao.

Sačekao je.

Crkva je ponovo rasla.

Kruciforma je postala temelj novog poretka.

Vaskrsenje je povezalo milijarde ljudi sa sistemom koji nisu razumeli.

A negde unutar svega toga Srž je ponovo pronašla svoje mesto.

\bigskip

Rol je razmišljao o kruciformama.

O ljudima koji mogu umreti tokom putovanja i vratiti se po dolasku.

O vojsci koja smrt više ne smatra konačnom.

O Crkvi koja kontroliše taj dar.

\bigskip

Odjednom mu je stara odluka da nikada ne primi kruciformu izgledala važnije nego ikada.

\bigskip

Anea je rekla još nešto.

TehnoCentar ne otvara portale pred njima.

\bigskip

``Ako ih Srž ne otvara, ko ih otvara?''

``Ne znam.''

``Hoćemo li saznati?''

Anea je pogledala prema starom luku kojim treba da prođu.

``Ako preživimo.''

\bigskip

Sledeći svet bio je Božji Gaj.

Anea je znala da ih tamo nešto čeka.

Ne De Soja.

Ne obični vojnici Paksa.

Nešto drugo.

Nešto iz Srži.

\bigskip

Rol je predložio da ne prođu kroz portal.

Mogli su ostati na Kom-Rijadu.

Sakriti se.

Promeniti put.

\bigskip

Anea je odmah odmahnula glavom.

``Moramo dalje.''

``Zašto?''

``Zato što se tamo put nastavlja.''

``To nije dovoljan razlog da uđemo u zasedu.''

``Za mene jeste.''

\bigskip

Rol je osetio stari bes.

Martin mu je rekao da je zaštiti.

Kako da je zaštiti od odluke koju sama donosi?

\bigskip

``Ja mogu da te zaustavim.''

Anea ga je pogledala.

``Možeš.''

``Mogu da te odnesem odavde.''

``Možeš.''

``I?''

``Nećeš.''

\bigskip

Rol je znao da je u pravu.

Ako je silom spreči da izabere sopstveni put, onda više nije njen zaštitnik.

Postao bi samo druga vrsta tamničara.

\bigskip

A. Betik je stajao između njih.

``Č. Endimione'', rekao je tiho, ``verujem da smo svi znali da će ovaj trenutak doći.''

Rol ga je pogledao.

``Naravno da si ti na njenoj strani.''

Androidov izraz jedva se promenio.

``Mislio sam da više nemamo strane.''

\bigskip

Rol je uzdahnuo.

``Mrzim kada si razuman.''

``To mi je jedna od malobrojnih zabava, ser.''

\bigskip

Pred zoru su krenuli prema splavu.

Šrajk je još bio tamo.

Čekao ih je kraj vode.

\bigskip

Rol je ponovo posegnuo za oružjem.

Anea je samo prošla pored metalnog stvorenja.

A. Betik za njom.

Rol je ostao poslednji.

Crvene Šrajkove oči pratile su ga.

\bigskip

``Ako je povrediš'', rekao je Rol tiho, ``znam da ovo neće imati nikakvog smisla, ali pokušaću da te ubijem.''

Šrajk se nije pomerio.

\bigskip

Rol je spustio oružje i pridružio se ostalima.

\bigskip

Splav je krenuo niz reku.

Iza njih su ostajali Hiperion, Parvati, Renesansa Vektor, Mare Infinitus, prazni Hebron, zaleđeni Sol Drakoni Septem i Kom-Rijad.

Svaki svet dao im je deo odgovora.

Ali sa svakim odgovorom broj pitanja je rastao.

\bigskip

Zašto se mrtvi dalekobacači otvaraju pred Aneom?

Ko ih kontroliše?

Šta se dogodilo stanovnicima Hebrona?

Zašto je TehnoCentar ponovo povezan sa Paksom?

Šta Šrajk želi?

I zašto je nešto iz Srži spremno da ubija sve koji se nađu na Aneinom putu?

\bigskip

Rol više nije verovao da samo beže.

Neko ih je vodio.

A neko drugi pokušavao da ih zaustavi.

\bigskip

Pred njima je bio Božji Gaj.

I prvi protivnik koji nije nameravao da Aneu zarobi.

Nameravao je da je ubije.